\chapter{振動・波動}
\section{局所的な振動から伝播へ}
振動は自由度が時間的に変わる運動であり、波動は隣接する自由度の結合を通して擾乱が空間へ伝わる現象である。
ばねの単振動と小角の振り子は第\ref{theory:mechanics:oscillation}節で扱った。
ここでは、変位の運動と位相の伝播を区別する。波が右へ進んでも媒質の各粒子が右へ運ばれ続けるとは限らない。
地震探査、音響測定、信号伝送へ共通するのは、媒質の復元力が速さを決め、幾何と境界が位相差を決めることである。

\section{位相を空間と時間で読む}\label{theory:wave:phase}
正弦波を \(y=A\sin(\omega t-kx+\phi_0)\) とする。
\(\omega\) は時間に対する位相の変化率、\(k\) は空間に対する位相の変化率であり、単位はそれぞれ rad/s、rad/m である。
時間又は空間で位相が \(2\pi\) 変化する間隔を定めれば
\[
 T=\frac{2\pi}{\omega},\qquad \lambda=\frac{2\pi}{k},\qquad f=\frac1T
\]
となる。同じ山を追うには \(\omega\,dt-k\,dx=0\) を満たすので、位相速度は
\[
 v_{\mathrm{ph}}=\frac{dx}{dt}=\frac{\omega}{k}=\frac{\lambda}{T}
\]
である。\(kx+\omega t\) と書けば進行方向は逆になる。
この結果は位相の運動を表し、分散性の媒質では波束の群速度 \(d\omega/dk\) と同じとは限らない。
試験の単一正弦波へ、波束を明示しないまま群速度を追加する必要はない。
周期で表した式の時間係数が \(2\pi/T\) であることと、位相速度に余分な \(2\pi\) が入らないことを一緒に確認する。

\section{復元力・密度と縦波・横波}\label{theory:wave:elastic}
等方線形弾性体の変位 \(\boldsymbol u\) に対する運動方程式を、密度 \(\rho\)、体積弾性率 \(K\)、剛性率 \(\mu\) で書くと
\[
 \rho\,\partial_t^2\boldsymbol u
 =(K+\mu/3)\nabla(\nabla\cdot\boldsymbol u)+\mu\nabla^2\boldsymbol u
\]
である。波数方向に平行な変位と垂直な変位へ分解することで
\[
 c_P=\sqrt{\frac{K+4\mu/3}{\rho}},\qquad
 c_S=\sqrt{\frac{\mu}{\rho}}
\]
を得る。縦波には体積変形とせん断変形の双方の復元力が入り、横波にはせん断の復元力が入る。
安定な等方モデルでは \(K,\mu>0\) であり、\(c_P>c_S\) となる。
これは異方性結晶の任意の方向へ無条件に適用する式ではない。
通常の流体は静的せん断に抗する \(\mu\) を持たず、体積内の音は圧縮の縦波である。
真空では音を伝える媒質がない。液体表面の重力波などを、流体の体積音波と同じ分類にしない。

気体では音波の小振幅の圧縮・膨張をほぼ断熱とみなして \(K_s=\gamma p\) とする。
理想気体式 \(p=\rho RT/M\) より
\[
 c=\sqrt{\frac{\gamma p}{\rho}}=\sqrt{\frac{\gamma RT}{M}}
\]
となる。\(T\) は絶対温度、\(R\) は気体定数、\(M\) はモル質量、\(\gamma=C_p/C_v\) である。
気体組成などを同じとして温度を増すと速くなる。
水は気体より密度も大きいが、体積弾性率の増大が支配的で、海水の音速は通常空気より大きい。
代表値は空気約340 m/s、海水約1500 m/sであり、精密測定では温度・圧力・塩分を指定する\cite{os_sound}。

\section{音源の移動と観測者の移動}\label{theory:wave:doppler}
音速 \(c\) を一定とし、速度は媒質に対して測る。
音源周期 \(T_s\) の間に波面は \(cT_s\) 進む。
音源が観測者から速さ \(u_s\) で遠ざかると、観測者側の波面間隔は \((c+u_s)T_s\) へ広がる。
静止観測者に対し
\[
 f_{\mathrm{obs}}=\frac{c}{(c+u_s)T_s}=f_s\frac{c}{c+u_s}
\]
となる。音源が近付けば分母は \(c-u_s\) となる。
一方、静止音源へ観測者が速さ \(u_o\) で近付けば、変わるのは波長ではなく波面の通過速度であり
\(f_{\mathrm{obs}}=f_s(c+u_o)/c\) となる。
両方が動く場合は、この二つの効果を別々に組み合わせる\cite{os_doppler}。
音源と観測者の相対速度だけを一つの分母へ入れる操作は一般に正しくない。

\section{動く船と反射音の距離}\label{theory:wave:echo}
音の反射を使う距離測定では、送受信位置が同じか移動したかが重要である。
船が岸壁へ速さ \(u\) で近付くとする。発音時距離を \(D_0\)、戻るまでの時間を \(t\) とすれば、
音は行きに \(D_0\)、帰りに受信時の距離 \(D_1=D_0-ut\) を進む。
\[
 ct=D_0+D_1=2D_0-ut,\qquad
 D_0=\frac{(c+u)t}{2},\qquad D_1=\frac{(c-u)t}{2}
\]
となる。岸壁での反射時刻を途中で明示しても同じ結果となる。
この式の二つの距離を区別すれば、「汽笛を鳴らしたとき」と「反射音を聞いたとき」の取り違えを防げる。
単に \(ct/2\) とするのは \(u=0\) の場合である。流れや風による音速の方向差は、問題が与えたときに別途扱う。

\section{二つの音源の経路差と位相差}\label{theory:wave:interference}
同じ角振動数の線形波を重ねると、受信点での合成振幅の二乗は
\[
 A^2=A_1^2+A_2^2+2A_1A_2\cos\delta,\qquad
 \delta=k(r_2-r_1)+\delta_0
\]
となる。\(r_i\) は経路長、\(\delta_0\) は放射時の相対位相である。
極小は \(\delta=(2n+1)\pi\) に現れるが、\(A_1\ne A_2\) なら残る振幅は \(|A_1-A_2|\) で、必ずしも無音ではない\cite{os_superposition}。
スピーカが間隔 \(a\) で並び、正面の聴取線まで \(L\) とする。
中央で強め合い、片方の正面で弱め合う設定なら、経路差は
\[
 \Delta=\sqrt{L^2+a^2}-L=\frac{a^2}{\sqrt{L^2+a^2}+L}
\]
である。後式は桁落ちを避けるにも有用である。
同位相の放射に対して \(\Delta=(2n+1)\lambda/2\) となり、
\[
 f_n=\frac{(2n+1)c}{2\Delta},\qquad n=0,1,2,\ldots
\]
である。最初の極小の指定がない場合、選択肢の周波数範囲も照合して次数を決める。
距離 \(L\) とスピーカ間隔 \(a\) の単純な差は経路差ではない。

\section{振動の合成・衝撃波・信号変調}\label{theory:wave:applications}
垂直な単振動 \(x=A\sin(\omega_1t+\alpha)\)、\(y=B\sin(\omega_2t+\beta)\) の軌道はリサジュー図形となる。
周波数比が有理なら周期的な閉軌道となり、位相差により直線や楕円などへ変わる。
一方、衝撃波は大振幅・非線形な圧縮で密度・圧力等が急に変わる前面である。
前面速度は未擾乱側の小振幅音速を超える。線形の正弦音波と同じ一定振幅モデルをそのまま当てない。
振幅変調では、例えば \([A_0+a\cos(\Omega t)]\cos(\omega_ct)\) と書き、
搬送波の振幅包絡へ信号を入れる。変調周波数 \(\Omega\) と搬送周波数 \(\omega_c\) を分ける。
これらは「波」という名称が共通でも、軌道・伝播前面・情報の表現という違う対象を問う。

\section{典型的な解法と出題の確認範囲}
正弦波では位相係数と次元、ドップラーでは媒質基準の速度、反射音では距離の時点、
干渉では経路差と次数を先に決める。
本章の振動3問に続き、S7で波動10問を原典確認した。
同じ正弦波や音源運動が再出題されても、選択肢の配列は一致しないことがある。
頻度の全体集計は索引・統計を参照し、この確認範囲だけから将来出題を断定しない。

\section{登録問題と解説初稿}
\section{音源運動・ドップラー・反射音}

\subsection{平成23年度 問3}
\label{q:h23-reg-applied-03}

\noindent\textbf{原典：} sources/H23/17\_applied\_science.pdf、PDF物理1頁。

\textbf{公式正答：} 2（single）。

\paragraph{問題の要旨} 静止した観測者から遠ざかる音源の観測振動数を求める。

\Needspace{4\baselineskip}\paragraph{確認した条件・記述（独自要約）}\mbox{}\par

\noindent\textbf{音源：} 440 Hz、観測者から4 m/sで遠ざかる。\par

\noindent\textbf{音速：} 340 m/s。観測者は静止。\par

\noindent\textbf{方向：} 観測振動数に最も近い値。\par

\paragraph{問われる核心} 媒質基準の音源運動によるドップラー効果

\paragraph{必要知識} 媒質基準の音源運動によるドップラー効果。\(f_{\rm obs}=f_s c/(c+u)=440\times340/344\simeq434.9\) Hzで2。音源が1周期の間に遠ざかる距離を波長へ加える。観測者の移動による波面の通過速度の変化とは分ける。

\paragraph{考え方・正答を選ぶ理由（査読待ち初稿）} \(f_{\rm obs}=f_s c/(c+u)=440\times340/344\simeq434.9\) Hzで2。音源が1周期の間に遠ざかる距離を波長へ加える。観測者の移動による波面の通過速度の変化とは分ける。

\Needspace{4\baselineskip}\paragraph{選択肢の検討（査読待ち初稿）}\mbox{}\par

\noindent\textbf{1：不適切} 430 Hz。後退で下がる方向は正しいが、計算値約434.9 Hzには435の方が近い。\par

\noindent\textbf{2：最適} 435 Hz。波面間隔が広がるため約434.9 Hzとなり、この候補が最も近い。\par

\noindent\textbf{3：不適切} 440 Hz。運動による波長変化を無視した音源そのものの振動数。\par

\noindent\textbf{4：不適切} 445 Hz。遠ざかるときの変化を接近と取り違えた候補。\par

\noindent\textbf{5：不適切} 450 Hz。遠ざかる音源で観測周波数が増えることはない。\par

\paragraph{再利用する考え方} 媒質基準の音源運動によるドップラー効果を暗記だけでなく、位相・幾何・符号と成立条件から導く。選択肢番号を類似年度へそのまま移さない。

\paragraph{共通理論} \hyperref[theory:wave:doppler]{第\ref*{theory:wave:doppler}節}

\par 参考：\cite{os_doppler}。

\paragraph{関連問題} \hyperref[q:h25-reg-applied-05]{H25-REG-APPLIED-05}

\par 難易度 2（校準初稿）、優先度未評価。

\subsection{平成25年度 問5}
\label{q:h25-reg-applied-05}

\noindent\textbf{原典：} sources/H25/17\_applied\_science.pdf、PDF物理3頁。

\textbf{公式正答：} 2（single）。

\paragraph{問題の要旨} 静止した観測者から遠ざかる音源の観測振動数を求める。

\Needspace{4\baselineskip}\paragraph{確認した条件・記述（独自要約）}\mbox{}\par

\noindent\textbf{音源：} 440 Hz、観測者から4 m/sで遠ざかる。\par

\noindent\textbf{音速：} 340 m/s。観測者は静止。\par

\noindent\textbf{方向：} 観測振動数に最も近い値。\par

\paragraph{問われる核心} 媒質基準の音源運動によるドップラー効果

\paragraph{必要知識} 媒質基準の音源運動によるドップラー効果。\(f_{\rm obs}=f_s c/(c+u)=440\times340/344\simeq434.9\) Hzで2。音源が1周期の間に遠ざかる距離を波長へ加える。観測者の移動による波面の通過速度の変化とは分ける。

\paragraph{考え方・正答を選ぶ理由（査読待ち初稿）} \(f_{\rm obs}=f_s c/(c+u)=440\times340/344\simeq434.9\) Hzで2。音源が1周期の間に遠ざかる距離を波長へ加える。観測者の移動による波面の通過速度の変化とは分ける。

\Needspace{4\baselineskip}\paragraph{選択肢の検討（査読待ち初稿）}\mbox{}\par

\noindent\textbf{1：不適切} 430 Hz。後退で下がる方向は正しいが、計算値約434.9 Hzには435の方が近い。\par

\noindent\textbf{2：最適} 435 Hz。波面間隔が広がるため約434.9 Hzとなり、この候補が最も近い。\par

\noindent\textbf{3：不適切} 440 Hz。運動による波長変化を無視した音源そのものの振動数。\par

\noindent\textbf{4：不適切} 445 Hz。遠ざかるときの変化を接近と取り違えた候補。\par

\noindent\textbf{5：不適切} 450 Hz。遠ざかる音源で観測周波数が増えることはない。\par

\paragraph{再利用する考え方} 媒質基準の音源運動によるドップラー効果を暗記だけでなく、位相・幾何・符号と成立条件から導く。選択肢番号を類似年度へそのまま移さない。

\paragraph{共通理論} \hyperref[theory:wave:doppler]{第\ref*{theory:wave:doppler}節}

\par 参考：\cite{os_doppler}。

\paragraph{関連問題} \hyperref[q:h23-reg-applied-03]{H23-REG-APPLIED-03}

\par 難易度 2（校準初稿）、優先度未評価。

\subsection{平成30年度 問5}
\label{q:h30-reg-applied-05}

\noindent\textbf{原典：} sources/H30/17\_applied\_science.pdf、PDF物理2頁。

\textbf{公式正答：} 2（single）。

\paragraph{問題の要旨} 岸壁へ接近する船が汽笛の反射音を聞いた時点の岸壁までの距離を求める。

\Needspace{4\baselineskip}\paragraph{確認した条件・記述（独自要約）}\mbox{}\par

\noindent\textbf{船速：} 岸壁へ20 m/sで等速接近。\par

\noindent\textbf{音速・時間：} 340 m/s、汽笛から反射音到着まで3 s。\par

\noindent\textbf{問う時点：} 反射音を聞いた時点。\par

\paragraph{問われる核心} 移動する観測位置と反射音の往復幾何

\paragraph{必要知識} 移動する観測位置と反射音の往復幾何。発音時の距離をD0、戻るまでt=\allowbreak{}3 sとすると音の総行程は\(ct=D_0+(D_0-ut)\)。\(D_0=(c+u)t/2=540\) m、戻った時点では\(D_1=D_0-ut=480\) mとなり2。

\paragraph{考え方・正答を選ぶ理由（査読待ち初稿）} 発音時の距離をD0、戻るまでt=\allowbreak{}3 sとすると音の総行程は\(ct=D_0+(D_0-ut)\)。\(D_0=(c+u)t/2=540\) m、戻った時点では\(D_1=D_0-ut=480\) mとなり2。

\Needspace{4\baselineskip}\paragraph{選択肢の検討（査読待ち初稿）}\mbox{}\par

\noindent\textbf{1：不適切} 450 m。船の移動60 mを正しい往復の幾何から減じた480 mとは一致しない。\par

\noindent\textbf{2：最適} 480 m。発音時540 mから移動60 mを差し引いた到着時距離。\par

\noindent\textbf{3：不適切} 510 m。静止する船としてct/2を計算した値で、船の移動を無視している。\par

\noindent\textbf{4：不適切} 540 m。発音時の距離であり、反射音を聞いた時点の距離ではない。\par

\noindent\textbf{5：不適切} 570 m。移動方向が接近であることと、音の往復距離の減少を正しく反映していない。\par

\paragraph{再利用する考え方} 移動する観測位置と反射音の往復幾何を暗記だけでなく、位相・幾何・符号と成立条件から導く。選択肢番号を類似年度へそのまま移さない。

\paragraph{共通理論} \hyperref[theory:wave:echo]{第\ref*{theory:wave:echo}節}

\par 難易度 2（校準初稿）、優先度未評価。

\section{音の経路差・位相差・干渉}

\subsection{令和元年度 問5}
\label{q:r01-reg-applied-05}

\noindent\textbf{原典：} sources/R01/17\_applied\_science.pdf、PDF物理2頁。

\textbf{公式正答：} 2（single）。

\paragraph{問題の要旨} 2スピーカの正面での強め合いと、片方の正面での極小から音の振動数を求める。

\Needspace{4\baselineskip}\paragraph{確認した条件・記述（独自要約）}\mbox{}\par

\noindent\textbf{配置：} スピーカ間隔2 m、正面の聴取線まで3.3 m。\par

\noindent\textbf{観測：} 中央は強め合い、片方の正面で最小。音速340 m/s。\par

\noindent\textbf{次数：} 初回極小という明示なし。候補範囲も照合。\par

\paragraph{問われる核心} 幾何学的経路差とコヒーレントな音の位相差

\paragraph{必要知識} 幾何学的経路差とコヒーレントな音の位相差。中央では両経路が等しく強め合う。片方の正面では距離\(L\)と\(\sqrt{L^2+a^2}\)の差\(\Delta=\sqrt{L^2+a^2}-L\)がある。逆位相の条件\(\Delta=(2n+1)\lambda/2\)の最初n=\allowbreak{}0では\(f=c/(2\Delta)\)≈304.2 Hzで2。次のn=\allowbreak{}1は約912.7 Hzで全候補より大きい。原典に最初の極小とは明示されないが、候補の範囲で次数を確認できる。距離差により振幅が異なるなら極小は必ずしも無音ではない。

\paragraph{考え方・正答を選ぶ理由（査読待ち初稿）} 中央では両経路が等しく強め合う。片方の正面では距離\(L\)と\(\sqrt{L^2+a^2}\)の差\(\Delta=\sqrt{L^2+a^2}-L\)がある。逆位相の条件\(\Delta=(2n+1)\lambda/2\)の最初n=\allowbreak{}0では\(f=c/(2\Delta)\)≈304.2 Hzで2。次のn=\allowbreak{}1は約912.7 Hzで全候補より大きい。原典に最初の極小とは明示されないが、候補の範囲で次数を確認できる。距離差により振幅が異なるなら極小は必ずしも無音ではない。

\Needspace{4\baselineskip}\paragraph{選択肢の検討（査読待ち初稿）}\mbox{}\par

\noindent\textbf{1：不適切} 280 Hz。最初の逆位相による極小の振動数約304.2 Hzから、正答候補より離れている。次の奇数次数は候補範囲外。\par

\noindent\textbf{2：最適} 310 Hz。道のり差約0.5588 mが半波長となる振動数約304.2 Hzに最も近い。\par

\noindent\textbf{3：不適切} 370 Hz。最初の逆位相による極小の振動数約304.2 Hzから、正答候補より離れている。次の奇数次数は候補範囲外。\par

\noindent\textbf{4：不適切} 430 Hz。最初の逆位相による極小の振動数約304.2 Hzから、正答候補より離れている。次の奇数次数は候補範囲外。\par

\noindent\textbf{5：不適切} 490 Hz。最初の逆位相による極小の振動数約304.2 Hzから、正答候補より離れている。次の奇数次数は候補範囲外。\par

\paragraph{再利用する考え方} 幾何学的経路差とコヒーレントな音の位相差を暗記だけでなく、位相・幾何・符号と成立条件から導く。選択肢番号を類似年度へそのまま移さない。

\paragraph{共通理論} \hyperref[theory:wave:interference]{第\ref*{theory:wave:interference}節}

\par 参考：\cite{os_superposition}。

\paragraph{関連問題} \hyperref[q:r05-reg-applied-05]{R05-REG-APPLIED-05}

\par 難易度 2（校準初稿）、優先度未評価。

\subsection{令和5年度 問5}
\label{q:r05-reg-applied-05}

\noindent\textbf{原典：} sources/R05/17\_applied\_science.pdf、PDF物理2頁。

\textbf{公式正答：} 3（single）。

\paragraph{問題の要旨} 2スピーカの正面での強め合いと、片方の正面での極小から音の振動数を求める。

\Needspace{4\baselineskip}\paragraph{確認した条件・記述（独自要約）}\mbox{}\par

\noindent\textbf{配置：} スピーカ間隔2.4 m、正面の聴取線まで4.0 m。\par

\noindent\textbf{観測：} 中央は強め合い、片方の正面で最小。音速340 m/s。\par

\noindent\textbf{次数：} 初回極小という明示なし。候補範囲も照合。\par

\paragraph{問われる核心} 幾何学的経路差とコヒーレントな音の位相差

\paragraph{必要知識} 幾何学的経路差とコヒーレントな音の位相差。中央では両経路が等しく強め合う。片方の正面では距離\(L\)と\(\sqrt{L^2+a^2}\)の差\(\Delta=\sqrt{L^2+a^2}-L\)がある。逆位相の条件\(\Delta=(2n+1)\lambda/2\)の最初n=\allowbreak{}0では\(f=c/(2\Delta)\)≈255.7 Hzで3。次のn=\allowbreak{}1は約767.2 Hzで全候補より大きい。原典に最初の極小とは明示されないが、候補の範囲で次数を確認できる。距離差により振幅が異なるなら極小は必ずしも無音ではない。

\paragraph{考え方・正答を選ぶ理由（査読待ち初稿）} 中央では両経路が等しく強め合う。片方の正面では距離\(L\)と\(\sqrt{L^2+a^2}\)の差\(\Delta=\sqrt{L^2+a^2}-L\)がある。逆位相の条件\(\Delta=(2n+1)\lambda/2\)の最初n=\allowbreak{}0では\(f=c/(2\Delta)\)≈255.7 Hzで3。次のn=\allowbreak{}1は約767.2 Hzで全候補より大きい。原典に最初の極小とは明示されないが、候補の範囲で次数を確認できる。距離差により振幅が異なるなら極小は必ずしも無音ではない。

\Needspace{4\baselineskip}\paragraph{選択肢の検討（査読待ち初稿）}\mbox{}\par

\noindent\textbf{1：不適切} 160 Hz。最初の逆位相による極小の振動数約255.7 Hzから、正答候補より離れている。次の奇数次数は候補範囲外。\par

\noindent\textbf{2：不適切} 210 Hz。最初の逆位相による極小の振動数約255.7 Hzから、正答候補より離れている。次の奇数次数は候補範囲外。\par

\noindent\textbf{3：最適} 260 Hz。道のり差約0.6648 mが半波長となる振動数約255.7 Hzに最も近い。\par

\noindent\textbf{4：不適切} 320 Hz。最初の逆位相による極小の振動数約255.7 Hzから、正答候補より離れている。次の奇数次数は候補範囲外。\par

\noindent\textbf{5：不適切} 400 Hz。最初の逆位相による極小の振動数約255.7 Hzから、正答候補より離れている。次の奇数次数は候補範囲外。\par

\paragraph{再利用する考え方} 幾何学的経路差とコヒーレントな音の位相差を暗記だけでなく、位相・幾何・符号と成立条件から導く。選択肢番号を類似年度へそのまま移さない。

\paragraph{共通理論} \hyperref[theory:wave:interference]{第\ref*{theory:wave:interference}節}

\par 参考：\cite{os_superposition}。

\paragraph{関連問題} \hyperref[q:r01-reg-applied-05]{R01-REG-APPLIED-05}

\par 難易度 2（校準初稿）、優先度未評価。

\section{音波・弾性波・媒質の復元力}

\subsection{平成24年度 問2}
\label{q:h24-reg-applied-02}

\noindent\textbf{原典：} sources/H24/17\_applied\_science.pdf、PDF物理1頁。

\textbf{公式正答：} 4（single）。

\paragraph{問題の要旨} 音の伝播、温度依存、流体の振動方向と固体の波速を判定する。

\Needspace{4\baselineskip}\paragraph{確認した条件・記述（独自要約）}\mbox{}\par

\noindent\textbf{方向：} 最も不適切。\par

\noindent\textbf{範囲：} 通常の音波。固体の比較は一般的な等方弾性体。\par

\paragraph{問われる核心} 圧縮波・せん断波の復元力と波速

\paragraph{必要知識} 圧縮波・せん断波の復元力と波速。媒質の体積弾性と密度が音速を決める。海水中の代表的な値は空気中より大きく、4が不適切。他の肢は通常の流体音波又は等方弾性体の一般的な説明として適切である。

\paragraph{考え方・正答を選ぶ理由（査読待ち初稿）} 媒質の体積弾性と密度が音速を決める。海水中の代表的な値は空気中より大きく、4が不適切。他の肢は通常の流体音波又は等方弾性体の一般的な説明として適切である。

\Needspace{4\baselineskip}\paragraph{選択肢の検討（査読待ち初稿）}\mbox{}\par

\noindent\textbf{1：適切} 真空中では音は伝わらない。。音は媒質の圧縮・変位の伝播であり、真空にはその媒質がない。\par

\noindent\textbf{2：適切} 空気の温度が高いほど音速は大きい。。理想気体の\(c=\sqrt{\gamma RT/M}\)では絶対温度の平方根に比例する。組成等を同じとして比較する。\par

\noindent\textbf{3：適切} 流体中を伝わる音は縦波である。。静的せん断弾性を持たない流体では、通常の体積音波は圧縮の縦波となる。表面波とは区別。\par

\noindent\textbf{4：不適切} 海水中の音速は空気中より小さい。。海水では概ね1500 m/s、空気では概ね340 m/sであり比較方向が逆。状態により変わるがこの大小は通常逆転しない。\par

\noindent\textbf{5：適切（一般的モデル）} 一般に固体中の縦波は横波より速い。。安定な等方線形弾性体で\(c_P^2-c_S^2=(K+\mu/3)/\rho>0\)である。異方性媒質の全方向をこの式だけで扱わない。\par

\paragraph{再利用する考え方} 圧縮波・せん断波の復元力と波速を暗記だけでなく、位相・幾何・符号と成立条件から導く。選択肢番号を類似年度へそのまま移さない。

\paragraph{共通理論} \hyperref[theory:wave:elastic]{第\ref*{theory:wave:elastic}節}

\par 参考：\cite{os_sound}。

\paragraph{関連問題} \hyperref[q:r04-reg-applied-05]{R04-REG-APPLIED-05}

\par 難易度 2（校準初稿）、優先度未評価。

\subsection{令和4年度 問5}
\label{q:r04-reg-applied-05}

\noindent\textbf{原典：} sources/R04/17\_applied\_science.pdf、PDF物理2頁。

\textbf{公式正答：} 2（single）。

\paragraph{問題の要旨} リサジュー図形、波の振動方向、衝撃波、縦横波の速さ、AMを判定する。

\Needspace{4\baselineskip}\paragraph{確認した条件・記述（独自要約）}\mbox{}\par

\noindent\textbf{方向：} 最も不適切。\par

\noindent\textbf{範囲：} 振動・波動の一般的説明。\par

\paragraph{問われる核心} 波の偏り、非線形な圧縮波、変調と単振動合成

\paragraph{必要知識} 波の偏り、非線形な圧縮波、変調と単振動合成。縦波・横波は速さではなく、変位方向と伝播方向の関係で定義する。同方向を横波とした2が不適切。固体の波速の一般的な比較は等方弾性体モデルで確認する。

\paragraph{考え方・正答を選ぶ理由（査読待ち初稿）} 縦波・横波は速さではなく、変位方向と伝播方向の関係で定義する。同方向を横波とした2が不適切。固体の波速の一般的な比較は等方弾性体モデルで確認する。

\Needspace{4\baselineskip}\paragraph{選択肢の検討（査読待ち初稿）}\mbox{}\par

\noindent\textbf{1：適切} 垂直な単振動の合成軌道は一般にリサジュー図形という。。xとyの単振動の周波数比・位相差で軌道が変わる。比が有理なら閉曲線となる。\par

\noindent\textbf{2：不適切} 伝播方向と媒質の振動方向が一致する波を横波という。。一致する波は縦波である。横波は伝播に垂直な振動を持つ。\par

\noindent\textbf{3：適切} 衝撃波は媒質中を超音速で伝播する圧縮波である。。不連続な圧力等の変化を伴う圧縮の前面が、未擾乱側の小振幅音速を超えて進む。小振幅正弦波とは分ける。\par

\noindent\textbf{4：適切（一般的モデル）} 一般に固体中の縦波は横波より速い。。安定な等方線形弾性体では縦波の復元力はK+4µ/3、横波はµなので、密度が同じなら縦波が速い。\par

\noindent\textbf{5：適切} AMは信号に応じて搬送波の振幅を変える方式である。。振幅包絡へ信号を入れる方式で、周波数を変えるFMと異なる。\par

\paragraph{再利用する考え方} 波の偏り、非線形な圧縮波、変調と単振動合成を暗記だけでなく、位相・幾何・符号と成立条件から導く。選択肢番号を類似年度へそのまま移さない。

\paragraph{共通理論} \hyperref[theory:wave:elastic]{第\ref*{theory:wave:elastic}節}, \hyperref[theory:wave:applications]{第\ref*{theory:wave:applications}節}

\par 参考：\cite{os_sound}。

\paragraph{関連問題} \hyperref[q:h24-reg-applied-02]{H24-REG-APPLIED-02}

\par 難易度 2（校準初稿）、優先度未評価。

\section{ばね・振り子}

\subsection{平成27年度 問1}
\label{q:h27-reg-applied-01}

\noindent\textbf{原典：} sources/H27/17\_applied\_science.pdf、PDF物理1頁。

\textbf{公式正答：} 4（single）。

\paragraph{問題の要旨} ばね定数1.0 N/mの軽いばねに質量8.0 kgをつるした鉛直単振動の周期を求める。

\Needspace{4\baselineskip}\paragraph{確認した条件・記述（独自要約）}\mbox{}\par

\noindent\textbf{spring\_N\_m：} 1.0\par

\noindent\textbf{mass\_kg：} 8.0\par

\noindent\textbf{spring\_mass：} negligible\par

\paragraph{問われる核心} 復元力の線形化と固有角振動数

\paragraph{必要知識} \(m\ddot x+kx=0\)なら\(T=2\pi\sqrt{m/k}\)。単振り子は小角近似で\(\ddot\theta+(g/\ell)\theta=0\)。

\paragraph{考え方（査読待ち初稿）} \(T=2\pi\sqrt{m/k}=2\pi\sqrt8\simeq17.8\) sで④。重力は平衡位置をずらすが、平衡位置からの振動周期には入らない。

\paragraph{正答を選ぶ理由} 独立に得た条件適合値・組合せは選択肢4である。導出した周期に最も近い。角周波数と周期を区別して2πを含める。

\Needspace{4\baselineskip}\paragraph{選択肢の検討（査読待ち初稿）}\mbox{}\par

\noindent\textbf{1：不採用} 1.4 s。指定条件を振動周期式へ代入した値と一致しない。\par

\noindent\textbf{2：不採用} 2.8 s。指定条件を振動周期式へ代入した値と一致しない。\par

\noindent\textbf{3：不採用} 8.8 s。指定条件を振動周期式へ代入した値と一致しない。\par

\noindent\textbf{4：採用} 18 s。導出した周期に最も近い。角周波数と周期を区別して2πを含める。\par

\noindent\textbf{5：不採用} 26 s。指定条件を振動周期式へ代入した値と一致しない。\par

\paragraph{再利用する考え方} 復元力の線形化と固有角振動数を、条件と式の対応から説明できること。

\paragraph{共通理論} \hyperref[theory:mechanics:oscillation]{第\ref*{theory:mechanics:oscillation}節}

\par 難易度 1（校準初稿）、優先度未評価。

\subsection{平成29年度 問1}
\label{q:h29-reg-applied-01}

\noindent\textbf{原典：} sources/H29/17\_applied\_science.pdf、PDF物理1頁。

\textbf{公式正答：} 3（single）。

\paragraph{問題の要旨} 長さ0.1 mのひもに質量0.1 kgの小さいおもりをつけた小角単振り子の周期を求める。

\Needspace{4\baselineskip}\paragraph{確認した条件・記述（独自要約）}\mbox{}\par

\noindent\textbf{length\_m：} 0.1\par

\noindent\textbf{mass\_kg：} 0.1\par

\noindent\textbf{g\_m\_s2：} 9.8\par

\noindent\textbf{angle：} small\par

\noindent\textbf{bob\_size：} negligible\par

\paragraph{問われる核心} 復元力の線形化と固有角振動数

\paragraph{必要知識} \(m\ddot x+kx=0\)なら\(T=2\pi\sqrt{m/k}\)。単振り子は小角近似で\(\ddot\theta+(g/\ell)\theta=0\)。

\paragraph{考え方（査読待ち初稿）} 小角近似\(\sin\theta\simeq\theta\)から\(T=2\pi\sqrt{\ell/g}\)。\(T=2\pi\sqrt{0.1/9.8}\simeq0.635\) sで③。質量は消える。

\paragraph{正答を選ぶ理由} 独立に得た条件適合値・組合せは選択肢3である。導出した周期に最も近い。角周波数と周期を区別して2πを含める。

\Needspace{4\baselineskip}\paragraph{選択肢の検討（査読待ち初稿）}\mbox{}\par

\noindent\textbf{1：不採用} 0.1 s。指定条件を振動周期式へ代入した値と一致しない。\par

\noindent\textbf{2：不採用} 0.3 s。指定条件を振動周期式へ代入した値と一致しない。\par

\noindent\textbf{3：採用} 0.6 s。導出した周期に最も近い。角周波数と周期を区別して2πを含める。\par

\noindent\textbf{4：不採用} 1.0 s。指定条件を振動周期式へ代入した値と一致しない。\par

\noindent\textbf{5：不採用} 6.0 s。指定条件を振動周期式へ代入した値と一致しない。\par

\paragraph{再利用する考え方} 復元力の線形化と固有角振動数を、条件と式の対応から説明できること。

\paragraph{共通理論} \hyperref[theory:mechanics:oscillation]{第\ref*{theory:mechanics:oscillation}節}

\par 難易度 1（校準初稿）、優先度未評価。

\subsection{令和7年度 問2}
\label{q:r07-reg-applied-02}

\noindent\textbf{原典：} sources/R07/17\_applied\_science.pdf、PDF物理1頁。

\textbf{公式正答：} 1（single）。

\paragraph{問題の要旨} 質量0.80 kgを鉛直ばねにつるすと4.9 cm伸びた。平衡位置から引いて放す単振動の周期を求める。

\Needspace{4\baselineskip}\paragraph{確認した条件・記述（独自要約）}\mbox{}\par

\noindent\textbf{mass\_kg：} 0.8\par

\noindent\textbf{static\_extension\_m：} 0.049\par

\noindent\textbf{g\_m\_s2：} 9.8\par

\noindent\textbf{spring\_mass：} light\par

\paragraph{問われる核心} 復元力の線形化と固有角振動数

\paragraph{必要知識} \(m\ddot x+kx=0\)なら\(T=2\pi\sqrt{m/k}\)。単振り子は小角近似で\(\ddot\theta+(g/\ell)\theta=0\)。

\paragraph{考え方（査読待ち初稿）} 静的釣合い\(k\Delta\ell=mg\)から\(k=0.80\times9.8/0.049=160\) N/m。\(T=2\pi\sqrt{m/k}=2\pi\sqrt{0.80/160}\simeq0.444\) sで①。cmをmへ変換してから計算する。

\paragraph{正答を選ぶ理由} 独立に得た条件適合値・組合せは選択肢1である。導出した周期に最も近い。角周波数と周期を区別して2πを含める。

\Needspace{4\baselineskip}\paragraph{選択肢の検討（査読待ち初稿）}\mbox{}\par

\noindent\textbf{1：採用} \(4.4\times10^{-1}\) s。導出した周期に最も近い。角周波数と周期を区別して2πを含める。\par

\noindent\textbf{2：不採用} \(8.9\times10^{-1}\) s。指定条件を振動周期式へ代入した値と一致しない。\par

\noindent\textbf{3：不採用} 1.0 s。指定条件を振動周期式へ代入した値と一致しない。\par

\noindent\textbf{4：不採用} 1.8 s。指定条件を振動周期式へ代入した値と一致しない。\par

\noindent\textbf{5：不採用} 2.0 s。指定条件を振動周期式へ代入した値と一致しない。\par

\paragraph{再利用する考え方} 復元力の線形化と固有角振動数を、条件と式の対応から説明できること。

\paragraph{共通理論} \hyperref[theory:mechanics:oscillation]{第\ref*{theory:mechanics:oscillation}節}

\par 難易度 2（校準初稿）、優先度未評価。

\section{進行波の位相・波長・角振動数}

\subsection{平成25年度 問2}
\label{q:h25-reg-applied-02}

\noindent\textbf{原典：} sources/H25/17\_applied\_science.pdf、PDF物理1頁。

\textbf{公式正答：} 3（single）。

\paragraph{問題の要旨} 正弦進行波の位相から波長と位相速度の組合せを選ぶ。

\Needspace{4\baselineskip}\paragraph{確認した条件・記述（独自要約）}\mbox{}\par

\noindent\textbf{波：} \(y=A\sin(\omega t-kx)\)、A・ω・kは正。\par

\noindent\textbf{方向：} 波長λと伝播速度vの適切な組合せ。\par

\paragraph{問われる核心} 進行波の位相・波数・角振動数

\paragraph{必要知識} 進行波の位相・波数・角振動数。空間で同位相が繰り返す条件は\(k\lambda=2\pi\)。一定位相を微分すると\(\omega\,dt-k\,dx=0\)で\(v=dx/dt=\omega/k\)。両式が一致する3を選ぶ。角振動数と周期を取り違えない。

\paragraph{考え方・正答を選ぶ理由（査読待ち初稿）} 空間で同位相が繰り返す条件は\(k\lambda=2\pi\)。一定位相を微分すると\(\omega\,dt-k\,dx=0\)で\(v=dx/dt=\omega/k\)。両式が一致する3を選ぶ。角振動数と周期を取り違えない。

\Needspace{4\baselineskip}\paragraph{選択肢の検討（査読待ち初稿）}\mbox{}\par

\noindent\textbf{1：不適切} \(\lambda=2\pi/\omega,\ v=2\pi\omega/k\)。波長の分母は波数であり、角振動数ではない。速度に余分な2πもある。\par

\noindent\textbf{2：不適切} \(\lambda=2\pi/\omega,\ v=2\pi k/\omega\)。波長は角振動数の逆数ではない。速度は逆数に加えて2πが余分。\par

\noindent\textbf{3：最適} \(\lambda=2\pi/k,\ v=\omega/k\)。空間位相が2π増える距離と位相一定の移動から導かれ、両方正しい。\par

\noindent\textbf{4：不適切} \(\lambda=2\pi/k,\ v=k/\omega\)。k/ωは速度の逆数となり、速度の次元が違う。\par

\noindent\textbf{5：不適切} \(\lambda=2\pi/k,\ v=2\pi\omega/k\)。波長は正しいが位相速度に余分な2πを掛けている。\par

\paragraph{再利用する考え方} 進行波の位相・波数・角振動数を暗記だけでなく、位相・幾何・符号と成立条件から導く。選択肢番号を類似年度へそのまま移さない。

\paragraph{共通理論} \hyperref[theory:wave:phase]{第\ref*{theory:wave:phase}節}

\paragraph{関連問題} \hyperref[q:h29-reg-applied-04]{H29-REG-APPLIED-04}

\par 難易度 2（校準初稿）、優先度未評価。

\subsection{平成29年度 問4}
\label{q:h29-reg-applied-04}

\noindent\textbf{原典：} sources/H29/17\_applied\_science.pdf、PDF物理2頁。

\textbf{公式正答：} 5（single）。

\paragraph{問題の要旨} 正弦進行波の位相から波長と位相速度の組合せを選ぶ。

\Needspace{4\baselineskip}\paragraph{確認した条件・記述（独自要約）}\mbox{}\par

\noindent\textbf{波：} \(y=A\sin(\omega t-kx)\)、A・ω・kは正。\par

\noindent\textbf{方向：} 波長λと伝播速度vの適切な組合せ。\par

\paragraph{問われる核心} 進行波の位相・波数・角振動数

\paragraph{必要知識} 進行波の位相・波数・角振動数。空間で同位相が繰り返す条件は\(k\lambda=2\pi\)。一定位相を微分すると\(\omega\,dt-k\,dx=0\)で\(v=dx/dt=\omega/k\)。両式が一致する5を選ぶ。角振動数と周期を取り違えない。

\paragraph{考え方・正答を選ぶ理由（査読待ち初稿）} 空間で同位相が繰り返す条件は\(k\lambda=2\pi\)。一定位相を微分すると\(\omega\,dt-k\,dx=0\)で\(v=dx/dt=\omega/k\)。両式が一致する5を選ぶ。角振動数と周期を取り違えない。

\Needspace{4\baselineskip}\paragraph{選択肢の検討（査読待ち初稿）}\mbox{}\par

\noindent\textbf{1：不適切} \(\lambda=2\pi/\omega,\ v=2\pi\omega/k\)。波長の分母は波数であり、角振動数ではない。速度に余分な2πもある。\par

\noindent\textbf{2：不適切} \(\lambda=2\pi/\omega,\ v=2\pi k/\omega\)。波長は角振動数の逆数ではない。速度は逆数に加えて2πが余分。\par

\noindent\textbf{3：不適切} \(\lambda=2\pi/k,\ v=2\pi\omega/k\)。波長は正しいが位相速度に余分な2πを掛けている。\par

\noindent\textbf{4：不適切} \(\lambda=2\pi/k,\ v=k/\omega\)。k/ωは速度の逆数となり、速度の次元が違う。\par

\noindent\textbf{5：最適} \(\lambda=2\pi/k,\ v=\omega/k\)。空間位相が2π増える距離と位相一定の移動から導かれ、両方正しい。\par

\paragraph{再利用する考え方} 進行波の位相・波数・角振動数を暗記だけでなく、位相・幾何・符号と成立条件から導く。選択肢番号を類似年度へそのまま移さない。

\paragraph{共通理論} \hyperref[theory:wave:phase]{第\ref*{theory:wave:phase}節}

\paragraph{関連問題} \hyperref[q:h25-reg-applied-02]{H25-REG-APPLIED-02}, \hyperref[q:r07-reg-applied-05]{R07-REG-APPLIED-05}

\par 難易度 2（校準初稿）、優先度未評価。

\subsection{令和7年度 問5}
\label{q:r07-reg-applied-05}

\noindent\textbf{原典：} sources/R07/17\_applied\_science.pdf、PDF物理2頁。

\textbf{公式正答：} 4（single）。

\paragraph{問題の要旨} 周期Tと波長λで書かれた正弦波の位相速度と角振動数を求める。

\Needspace{4\baselineskip}\paragraph{確認した条件・記述（独自要約）}\mbox{}\par

\noindent\textbf{波：} \(y=A\sin\{2\pi(t/T-x/\lambda)\}\)。A,T,λは正。\par

\noindent\textbf{方向：} 伝播速度vと角振動数ωの適切な組合せ。\par

\noindent\textbf{候補5：} 角振動数は2πλT（原画像照合、OCR誤認を採用しない）。\par

\paragraph{問われる核心} 周期・波長による位相速度と角振動数

\paragraph{必要知識} 周期・波長による位相速度と角振動数。\(y=A\sin\{2\pi(t/T-x/\lambda)\}\)の時間係数は\(\omega=2\pi/T\)。一定位相では\(dt/T=dx/\lambda\)なので\(v=\lambda/T\)で4。

\paragraph{考え方・正答を選ぶ理由（査読待ち初稿）} \(y=A\sin\{2\pi(t/T-x/\lambda)\}\)の時間係数は\(\omega=2\pi/T\)。一定位相では\(dt/T=dx/\lambda\)なので\(v=\lambda/T\)で4。

\Needspace{4\baselineskip}\paragraph{選択肢の検討（査読待ち初稿）}\mbox{}\par

\noindent\textbf{1：不適切} \(v=2\pi/\lambda,\ \omega=\lambda/T\)。速度と角振動数の双方の次元が異なる。\par

\noindent\textbf{2：不適切} \(v=\lambda T,\ \omega=2\pi/\lambda\)。速度は周期で割る量であり積ではない。角振動数も空間の波数ではない。\par

\noindent\textbf{3：不適切} \(v=2\pi\lambda/T,\ \omega=2\pi/T\)。角振動数は正しいが、位相速度に余分な2πがある。\par

\noindent\textbf{4：最適} \(v=\lambda/T,\ \omega=2\pi/T\)。位相一定の微分と時間位相係数を両方満たす。\par

\noindent\textbf{5：不適切} \(v=\lambda/T,\ \omega=2\pi\lambda T\)。速度は正しいが角振動数に波長と周期を掛け、次元が異なる。\par

\paragraph{再利用する考え方} 周期・波長による位相速度と角振動数を暗記だけでなく、位相・幾何・符号と成立条件から導く。選択肢番号を類似年度へそのまま移さない。

\paragraph{共通理論} \hyperref[theory:wave:phase]{第\ref*{theory:wave:phase}節}

\paragraph{関連問題} \hyperref[q:h29-reg-applied-04]{H29-REG-APPLIED-04}

\par 難易度 2（校準初稿）、優先度未評価。

